\begin{table}[tbp]
\centering
\caption{\textbf{How the sample was built, and what was left out.} The upper
panel narrows the panel to program-years in which the same program can be
observed in two consecutive years with its major-program status known in both.
The lower panel is the reason the paper does not simply code a missing prior
year as a clean year. In 366,506 program-years --- 15.35 percent of the panel --- the
auditee filed an audit last year but this program does not appear in it. Those
program-years record a finding 3.26 percent of the time, which is not the rate of
a clean program. They are excluded rather than recoded, because the absence of a
record is not a clean result.}
\label{tab:sample}
\footnotesize
\setlength{\tabcolsep}{3pt}
\begin{tabular}{lrr}
\toprule
Step & Program-years & Excluded here \\
\midrule
All program-years in the panel & 2,387,511 &  \\
Within the \$750,000 regime (FY2016--2023) & 2,387,511 & 0 \\
With a resolved prior-year observation & 1,600,898 & 786,613 \\
Sample A: audited as a major program in both years & 206,926 & 1,393,972 \\
\bottomrule
\end{tabular}

\vspace{1.1em}

\footnotesize
\begin{tabular}{lrrrl}
\toprule
What the prior year looks like & Program-years & Share & \shortstack{Finding rate\\this year} & Treatment \\
\midrule
Prior year resolved, finding recorded & 66,131 & 2.77\% & 0.4723 & kept \\
Prior year resolved, no finding recorded & 1,534,767 & 64.28\% & 0.0206 & kept \\
Auditee filed last year; this program is not in that filing & 366,506 & 15.35\% & 0.0326 & excluded \\
No filing at all for this auditee last year & 160,861 & 6.74\% & 0.0706 & excluded \\
First year of the window (2016) & 259,246 & 10.86\% & 0.0462 & excluded \\
\bottomrule
\end{tabular}
\end{table}
